\nonstopmode
\documentclass[a4paper, 11pt]{article}
\usepackage{graphicx} % Required for inserting images
\usepackage{amsmath}
\usepackage{geometry}
\usepackage{hyperref}
\usepackage{setspace}
\usepackage{xcolor}
\usepackage{colortbl} 
\usepackage{tabularray}
\definecolor{darkgreen}{RGB}{18,94,40}
\definecolor{lightgreen}{RGB}{179,255,179}
\definecolor{moregreen}{RGB}{153,255,143}
 \geometry{
 a4paper,
 left=25mm,
 right=25mm,
 top=20mm,
 bottom=20mm,
 }

\setlength{\parskip}{1em}
\setlength{\parindent}{0pt}
\graphicspath{{../../../media/}}

\begin{document}

\begin{minipage}{0.35\linewidth}
    \includegraphics[width=\linewidth]{Logo_Università_Padova.svg.png}
\end{minipage}\hfil
\begin{minipage}{0.55\linewidth}
\textbf{Università degli Studi di Padova} \\
Laurea in Informatica \\
Corso di Ingegneria del Software \\
Anno Accademico 2023/2024
\end{minipage}

\vspace{5mm}

\begin{minipage}{0.35\linewidth}
    \includegraphics[width=\linewidth]{logo rotondo.jpg}
\end{minipage}\hfil
\begin{minipage}{0.55\linewidth}
\textbf{Gruppo:} SWEet16 \\
\textbf{Email:} 
\href{mailto:sweet16.unipd@gmail.com}{\nolinkurl{sweet16.unipd@gmail.com}}
\end{minipage}

\vspace{15mm}

\begin{center}
\begin{Huge}
        \textbf{Verbale Interno} \\
        \vspace{4mm}
        \textbf{21 Novembre 2023}
\end{Huge}

\vspace{20mm}

\begin{large}
\begin{spacing}{1.4}
\begin{tabular}{c c c}
   Redattori:  &  Alberto M. & \\
   Verificatori: & Bilal E. & \\
   Amministratore: &  Alberto C. & \\
   Destinatari: & T. Vardanega & R. Cardin \\  
   Versione: & 1.0.0 & 
\end{tabular}
\end{spacing}
\end{large}
\end{center}

\pagebreak


\begin{huge}
    \textbf{Registro delle modifiche}
\end{huge}
\vspace{5pt}

\begin{tblr}{
colspec={|X[1.5cm]|X[2cm]|X[2.5cm]|X[2.5cm]|X[5cm]|},
row{odd}={bg=white},
row{even}={bg=lightgray},
row{1}={bg=black,fg=white}
}
    Versione & Data & Autore & Verificatore & Descrizione \\
    \hline
    1.0.0 & 2023/11/22 & Bilal E. & & Approvazione per il rilascio \\
    \hline
     0.1.0 & 2023/11/21 & Alberto C. & Bilal E. & Prima stesura del documento \\
     \hline
\end{tblr}

\pagebreak

\section{Partecipanti}
Di seguito i nomi dei partecipanti con le rispettive matricole: \\
\vspace{5mm}

\begin{table}[h]
\begin{tblr}{
colspec={X[5cm]X[5cm]},
row{odd}={bg=moregreen},
row{even}={bg=lightgreen},
row{1}={bg=darkgreen,fg=white}
}
    Nome & Matricola \\
    Alberto Cinel & 1142833 \\
    Bilal El Moutaren & 2053470 \\
    Alberto Michelazzo & 2010007 \\
    Alex Scantamburlo & 2042326 \\
    Iulius Signorelli & 2012434 \\
    Giovanni Zuliani & 595900 
\end{tblr}
\end{table}

\vspace{10pt}

\textbf{Inizio incontro:} Ore 21:00 \newline
\textbf{Fine incontro:} Ore 22:30  \newline

\pagebreak

\section{Sintesi ed elaborazione incontro}

Nella prima parte dell'incontro si è discusso brevemente sul progresso dei documenti, ci si è concentrati in particolare sull'analisi dei requisiti. \\
Si è deciso di introdurre una milestone, con data di scadenza fissata al 05/12/2023, entro il quale tre componenti del gruppo si coadiuveranno per analizzare in maniera preliminare i vari casi d'uso; due altri membri si organizzeranno, invece, per inizializzare il PoC. \\
I restanti componenti si occuperanno della verifica del documento precedente citato.

Si è inoltre deciso di organizzare un secondo incontro con il sig. Staffolani per discutere dei requisiti richiesti per l'MVP e le funzionalità da implementare nel POC. \\
Sarà inoltre chiesto se si terranno corsi di formazione da parte dell'azienda sulle tecnologie implementative da noi non conosciute, in particolar modo su React come linguaggio di front-end e NodeJS come linguaggio di back-end.

Infine, dopo una rapida discussione sui vari linguaggi da poter utilizzare per la parte front-end, si sono distribuite le mansioni ai membri del gruppo. 

Ci si aggiornerà successivamente tramite Telegram per la data dell'incontro con con il referente di \textit{Imola Informatica}.

Si conclude l'incontro decidendo martedì 28/11 come data del prossimo ritrovo sincrono.

\end{document}



